\documentclass[11pt]{article}
\usepackage[margin=1in]{geometry}
\usepackage{amsmath,amssymb,amsthm,booktabs,enumitem}
\usepackage[hidelinks]{hyperref}

\newcommand{\bps}{\,\text{bps}}
\newcommand{\usd}{\$}
\title{Route Right: A Cost Model and Allocation Engine\\ for Multi-Venue Crypto Execution}
\author{Engine specification, v0.1}
\date{\today}

\begin{document}
\maketitle

\begin{abstract}
We formalize the per-period (trailing 30-day) cost a trading firm incurs when
its order flow is split across crypto venues with \emph{status-based tiered}
fee schedules, native-token discounts, conditional market-maker programs, and
finite liquidity. The central structural fact is that explicit fees are
non-convex in volume (tier rates step \emph{down} at thresholds, rewarding
concentration) while market impact is convex (penalizing it); the optimal route
balances the two. We give the cost function, prove the relevant monotonicities,
and reduce the allocation problem to a mixed-integer linear program solved
exactly at the scale of interest.
\end{abstract}

\section{Setup and notation}
Let $\mathcal V$ index venues. Over one rolling 30-day window a firm must
execute total notional $V>0$ in $N$ executions, so the average trade size is
$S = V/N$. A fraction $\varphi\in[0,1]$ of notional is executed as
\emph{taker} (crossing the spread); the complement $1-\varphi$ is
\emph{maker} (passive). The decision is an allocation
\[
x = (x_v)_{v\in\mathcal V}, \qquad x_v\ge 0, \qquad \sum_{v} x_v = V,
\]
where $x_v$ is notional routed to venue $v$. We assume the firm preserves its
order-size distribution across venues, so the executions placed on $v$ number
$n_v = N\,x_v/V$ and retain average size $S$.

Rates are stored as fractions of notional; $0.0010 = 0.10\% = 10\bps$. A
negative maker rate denotes a rebate.

\section{Explicit fees}
\subsection{Tiers as a status}
Each venue has an ascending tier ladder with thresholds
$0=\Theta_{v,0}<\Theta_{v,1}<\cdots<\Theta_{v,K_v}$ and per-tier maker/taker
rates $(m_{v,k},t_{v,k})$. Crucially, the tier is a \emph{status}: the rate of
the attained tier applies to the firm's \emph{entire} volume on $v$, not
marginally. The attained tier index is
\[
k_v(q_v) = \max\{k : q_v \ge \Theta_{v,k}\},
\]
where the qualifying metric $q_v$ depends on the venue's rule: volume
($q_v=x_v$), assets ($q_v=A_v$), \texttt{either} ($q_v=\max(x_v,A_v)$), or
\texttt{both} ($q_v=\min(x_v,A_v)$, e.g.\ a venue requiring both volume and a
token balance).

\subsection{Effective rates}
Two adjustments turn posted rates into expected effective rates.

\paragraph{Native-token discount.} If the firm pays fees in $v$'s token, a
discount $\delta_v$ multiplies \emph{positive} fees only (rebates are
unaffected), at the cost of a carry $c^{\text{tok}}_v$ (bps on routed notional)
for holding the token. The effective taker rate is
$\tilde t_{v,k} = t_{v,k}\,(1-\delta_v\,\mathbf 1[\text{pay in token}])$ when
$t_{v,k}>0$.

\paragraph{Expected maker rate.} A passive order earns its (possibly negative)
rate only if it \emph{fills}. With fill probability $p_v$, adverse-selection
loading $a_v$, and a chase cost $\approx$ a half-spread $h_v$ on the unfilled
remainder,
\begin{equation}
\tilde m_{v,k} \;=\; p_v\bigl(m^{\,\delta}_{v,k} + a_v\bigr) \;+\; (1-p_v)\,h_v,
\label{eq:maker}
\end{equation}
where $m^{\,\delta}_{v,k}$ is the token-adjusted maker rate. Equation
\eqref{eq:maker} is the correction that prevents the engine from treating a
posted rebate as guaranteed income.

\paragraph{Market-maker program (conditional).} If the firm is eligible and
$\Theta_{v,k}\ge \Theta^{\text{mm}}_v$, an enhanced pair
$(m^{\text{mm}}_v,t^{\text{mm}}_v)$ is available in exchange for an obligation
cost $o_v$ (bps on maker notional). The engine adopts the program on $v$ iff it
lowers the blended rate.

The blended explicit rate at tier $k$ is
\begin{equation}
r_{v,k} \;=\; \varphi\,\tilde t_{v,k} \;+\; (1-\varphi)\,\tilde m_{v,k}
\;+\; o_v\,\mathbf 1[\text{mm}] \;+\; c^{\text{tok}}_v\,\mathbf 1[\text{token}],
\label{eq:blend}
\end{equation}
and the explicit fee on $v$ is $r_{v,k_v}\,x_v$, a \emph{piecewise-linear,
non-convex} function of $x_v$ (it drops discontinuously in slope at each
$\Theta_{v,k}$).

\begin{proposition}[Tier monotonicity]
For fixed $\varphi$ and parameters, $r_{v,k}$ is non-increasing in $k$.
Consequently the blended rate is a non-increasing step function of the
qualifying metric, and average cost per dollar weakly decreases with volume on a
venue. This concavity is the economic incentive to concentrate flow.
\end{proposition}

\section{Implicit cost: spread and impact}
\paragraph{Spread.} Takers cross the spread each trade:
$\;\text{spread}_v(x_v) = \varphi\,x_v\,h_v$, linear in $x_v$.

\paragraph{Participation impact (convex).} Absorbing a share of a venue's window
volume $\mathrm{ADV}_v$ moves price by a square-root law, giving an aggregate
cost
\begin{equation}
I^{\text{part}}_v(x_v) \;=\; \kappa\,\sigma\,x_v\sqrt{\frac{x_v}{\mathrm{ADV}_v}}
\;=\; \kappa\,\sigma\,\mathrm{ADV}_v^{-1/2}\,x_v^{3/2},
\label{eq:part}
\end{equation}
\emph{convex} and increasing. This is the counterweight to tier concavity: it
makes spreading worthwhile and yields an interior optimum.

\paragraph{Size impact (trade-count effect).} Holding notional fixed, fewer
(larger) orders dig deeper into the instantaneous book. With instantaneous depth
$d_v=\beta\,\mathrm{ADV}_v$,
\begin{equation}
I^{\text{size}}_v(x_v) \;=\; \kappa_S\,\sigma\,x_v\sqrt{\frac{S}{d_v}},
\qquad S=\frac{V}{N},
\label{eq:size}
\end{equation}
linear in $x_v$ but scaling as $\sqrt{S}=\sqrt{V/N}$, so for fixed $V$ the cost
\emph{rises as $N$ falls}. Equations \eqref{eq:part}--\eqref{eq:size} are how the
engine honors the notional-versus-trade-count distinction: $V$ sets tiers, $S$
sets per-order impact. Optional flat per-order fees add $n_v\,\phi_v$ (typically
$0$ for crypto spot).

\section{The allocation problem}
Total expected monthly cost is
\begin{equation}
C(x) = \sum_{v}\Bigl[\, r_{v,k_v(x_v)}\,x_v
\;+\; \varphi h_v x_v
\;+\; I^{\text{part}}_v(x_v) + I^{\text{size}}_v(x_v)
\;+\; n_v\phi_v \Bigr].
\label{eq:total}
\end{equation}
We minimize $C$ over $\sum_v x_v = V$, the counterparty cap
$x_v \le \rho V$ (share $\rho\in(0,1]$), access masks $x_v=0$ for venues the
firm's region cannot reach, and optional semi-continuous minimum lots. Because
the first term is non-convex, this is not a convex program.

\subsection{MILP formulation}
We disaggregate each $x_v$ over its tier brackets with binaries
$z_{v,k}\in\{0,1\}$ and segment variables $y_{v,k}\ge 0$:
\begin{align}
\sum_k z_{v,k} &= 1, & x_v &= \sum_k y_{v,k},\\
\Theta_{v,k}\,z_{v,k} \;\le\; y_{v,k} &\;\le\; U_{v,k}\,z_{v,k},
& U_{v,k}&=\min(\Theta_{v,k+1},\rho V),
\end{align}
so exactly one bracket is active and its rate applies to the whole volume; the
explicit term linearizes to $\sum_k r_{v,k}\,y_{v,k}$. The convex impact
\eqref{eq:part} is represented by a piecewise-linear outer description with
increasing slopes $s_{v,j}$ on segments $\delta_{v,j}\in[0,w_{v,j}]$,
$x_v=\sum_j\delta_{v,j}$; convexity makes the LP fill low-slope segments first,
so no extra binaries are needed and the approximation is exact at the
breakpoints. The remaining terms are linear in $x_v$. The result is a MILP with
$\sum_v K_v$ binaries.

\begin{remark}[Scale and exactness]
At the scale of interest ($|\mathcal V|\!\sim\!6\text{--}15$ venues, $K_v\!\sim\!6\text{--}10$
tiers) the MILP has on the order of $10^2$ binaries and CBC solves it to
proven optimality in milliseconds. Any single-venue allocation is feasible, so
the optimum is provably no worse than the naive ``best single venue''
baseline; the engine's edge is exactly the gap.
\end{remark}

\section{Rolling-tier sprint}
Tiers track \emph{trailing} 30-day volume, so mid-window a firm can push extra
volume now to re-rate remaining flow. With current trailing volume $V^{\text{now}}_v$
at day $\tau$ of a 30-day window, run-rate $\mu_v=V^{\text{now}}_v/\tau$, and
remaining flow $R_v=\mu_v(30-\tau)$, crossing the next threshold
$\Theta'$ (rate $r\to r'$) is worthwhile iff
\begin{equation}
\underbrace{(r-r')\,R_v}_{\text{saving on remaining flow}}
\;>\;
\underbrace{r\,\bigl(\Theta'-V^{\text{now}}_v\bigr)}_{\text{sprint cost}} .
\end{equation}
The engine reports the breakeven and net benefit per venue.

\section{What is data versus structure}
The \emph{structure}---tier non-convexity, convex impact, the MILP, the
sprint breakeven---is the engine's defensible content. The following are
\textbf{parameters to be calibrated} from live feeds and the firm's own fills,
and the engine's accuracy is bounded by their quality:
\begin{itemize}[nosep]
\item tier thresholds and rates (live per venue; top tiers and MM-program terms
are often negotiated/NDA);
\item half-spreads $h_v$, window liquidity $\mathrm{ADV}_v$, depth fraction
$\beta$, and impact coefficients $\kappa,\kappa_S$ (from order-book and TCA
data);
\item maker fill probabilities $p_v$ and adverse-selection loadings $a_v$ (from
the firm's realized passive fills).
\end{itemize}
Shipping the engine with indicative defaults is fine for illustration; shipping
it against a customer's real fills is what makes a route trustworthy.

\end{document}
